% Source: PAPER_MASTER_PLAN sec 7 (limitations) + the mandatory disclosures
% in CLAUDE.md. Frame honestly; a reviewer must not find an undisclosed hole.

\subsection{Two independent limits}
It is important not to conflate the two limits this study exposes. The
\emph{navigation} limit (Section~\ref{sec:res-limit}) is a property of the
sensing task: at high ranging noise the position of an obstacle is
uncertain to a degree comparable to its size, and no policy can recover
information the sensor never captured; this bounds the achievable ceiling
and is unaffected by any defense. The \emph{security} limit --- the
single-frame camouflage recall collapse --- is a property of the
\emph{detector}, and it is the one our temporal rule removes. The two are
orthogonal: our contribution restores success up to the noise-limited
ceiling, not beyond it, and never depresses it.

\subsection{On the precision caveat}
The temporal detector's precision at $\sigma=0.6$ is $0.80$--$0.82$, below a
nominal $0.9$ target, which under a purely detection-theoretic reading would
be a weakness. We argue precision is the wrong figure of merit here. It was
only ever a proxy for the harm caused by false gating, and that harm is
measured directly by the no-harm column, which is statistically
indistinguishable from zero. The reconciliation is that the residual false
flags fall on ultra-stealthy camouflage buckets --- phantoms hugging real
obstacles so tightly that they barely protrude. These are simultaneously the
hardest to distinguish from honest noise \emph{and} nearly harmless, which
is exactly the stealth/harm bind seen from the defender's side: a
misclassification there costs almost nothing. Own-LiDAR fusion and the
routed heading further cushion any wrongly gated broadcast. Tightening the
threshold to raise precision sacrifices recall on genuinely harmful
phantoms, so we adopt the operating point that maximizes recovered success
rather than detector precision.

\subsection{The verifiability boundary}
Two independent observations in the literature bound what any
consistency-based defense can achieve, and both are recovered by our
results. First, trust models are blind wherever verification is impossible:
Hallyburton and Pajic report that fabrications injected into regions no
other agent observes cannot be detected by existing
methods~\cite{aerialtrust2025}. Our camouflage attack is the statistical
analogue of that spatial blind spot --- the verifier \emph{does} observe the
region, but ranging noise forces a tolerance wide enough to hide the lie
inside it. Second, an attacker that mostly reports truthfully evades
trust accumulation: TruPercept observed that unreliable agents were not
separated from honest ones by trust scores~\cite{trupercept2020}. Our
adaptive-attacker study makes this quantitative and then closes it: evasion
requires collapsing the fabrication onto real geometry, at which point it
blocks no new space. Time, rather than a stronger single-frame test, is what
converts the unverifiable into the verifiable here --- and what remains
unverifiable is, by the same bind, unharmful.

\subsection{Robustness to trust-poisoning attacks}
A recognized failure mode of consistency-based trust is that an adversary
can engineer disagreements to make the defense downweight or exclude
\emph{honest} agents~\cite{trustflip2026}.
Our evaluation confronts this directly by reporting no-harm under an
adaptive, defense-aware attacker and finding it near zero: the temporal
test does not provide a lever by which induced disagreement translates into
honest exclusion, because it thresholds a time-averaged bias that honest
neighbours, by construction, do not accumulate.

\subsection{The warm-up window}
By construction the temporal test issues no verdict for a
(receiver, neighbour, track) bucket until $K_{\min}=20$ samples have
accumulated; during those first sightings the swarm is protected only by
the single-frame robust path. In our episodes (1200 steps) this window is
negligible, but a short-lived encounter, a late-joining traitor, or a
mission with brief pairwise contact would face reduced temporal protection
--- shrinking $K_{\min}$ trades this window against verdict reliability.

\subsection{Assumptions and their consequences}
We state the modelling assumptions and their effect on the claims. (i)~The goal-direction input to the
policy is not a straight-line bearing to the goal but the gradient of a
\emph{global} shortest-path distance field, recomputed by Dijkstra's
algorithm over the full obstacle layout at the start of each episode; each
drone is thereby handed an optimal obstacle-routed heading that presupposes a
global map it could not sense on its own. We model this as an external
mission planner that supplies routed waypoints, leaving the learned policy
responsible for local control and collision avoidance --- a conventional
division of labour, but a privileged input we disclose plainly rather than
fold into the sensing model. Its effect on our claims is conservative:
because the routed heading already steers each drone around the true obstacle
field, it \emph{dampens} the influence of a fabricated obstacle, so every
attack-induced success drop we report is a lower bound on what a planner-free
agent would suffer. Retraining the policy to navigate from a raw goal
coordinate, without the routed heading, is the single largest avenue to a
stronger navigation claim and the foremost limitation of the present study. (ii)~Communication is modelled as
lossless and zero-latency within a $10$~m range; real radios drop and delay
messages, which would reduce the evidence available to the temporal test and
warrants study under realistic channels. (iii)~Trust is applied at the
\emph{neighbour} level: one sustained contradiction excludes a neighbour's
entire contribution, discarding its honest observations too. Obstacle-level
gating would be less wasteful and is natural future work. (iv)~The trust
rule is hand-designed, not learned; we regard this as a feature, since we
find no learned component is necessary for this threat class, but it is
specific to the fabricated-obstacle model studied here. (v)~The study is in
simulation, with ten drones, in two dimensions, and with circular obstacles;
the mechanism is geometry-agnostic in principle, but validation on hardware,
in three dimensions, and with richer obstacle shapes remains open.
(vi)~The temporal test assumes sensor noise is \emph{unbiased}: its
separation of honest from fabricated relies on honest offsets being
zero-mean, so a persistent per-agent measurement bias would not cancel and
could cause an honest neighbour to be flagged. A bias common to all agents
cancels harmlessly in the pairwise offset; only \emph{agent-specific}
systematic bias is a concern, and we quantify it directly by sweeping a
synthetic per-agent bias against the same noise range ($f{=}2$, camouflage,
500 maps). No-harm stays statistically indistinguishable from zero for a
miscalibration up to $0.2$~m at every noise level tested; beyond that the
own-sensing cushion is no longer sufficient, detection precision falls from
$0.9$--$1.0$ to $0.3$--$0.4$, and no-harm reaches $-11.4$ pp at bias $0.4$~m
and $-18.2$~pp at bias $0.6$~m (both at $\sigma{=}0$, the noise level at
which the fixed component of the tolerance offers the least cover). The
effect is partly self-masking at high noise --- the noise-adaptive tolerance
$\varepsilon_0+k_\sigma\sigma$ widens with $\sigma$ and incidentally absorbs
part of the bias too --- but does not vanish. Distinguishing such a bias
from fabrication is moreover possible in principle --- a miscalibrated
sensor shifts \emph{all} of a neighbour's reported obstacles by the same
amount, whereas a camouflage lie displaces only the fabricated one, so
estimating and subtracting a per-neighbour global offset before the residual
bias test would separate the two --- which we leave to future work. Learned
per-agent pose-offset correction of exactly this kind has been shown effective
against systematic bias in benign collaborative
perception~\cite{vadivelu2020}, so the residual concern is a matter of
composing that correction with our test, not an unaddressed vulnerability.
